\section{Diseño y desarrollo del sistema}

El propósito central de esta fase fue llevar a la práctica los fundamentos teóricos y metodológicos previamente analizados, convirtiéndolos en un sistema de gestión de eventos plenamente funcional. Para ello, se tomó como base la arquitectura Modelo–Vista–Controlador (MVC), que facilitó la organización en capas; la Programación Orientada a Objetos (POO), que aportó principios de reutilización y escalabilidad; y la metodología ágil Scrum, que orientó el proceso de construcción mediante una gestión iterativa y colaborativa. En conjunto, estas guías permitieron estructurar y coordinar el proyecto de manera ordenada, asegurando tanto la calidad técnica del software como la efectividad en la organización del trabajo del equipo.